\section{Билет 45. Первое начало термодинамики. Смысл и знаки входящих в него величин. Понятие количества теплоты. Первое начало термодинамики для изопроцессов.}

\begin{definition}
\textbf{Первое начало термодинамики} является обобщением закона сохранения энергии на тепловые процессы. Оно формулируется следующим образом:
\begin{equation}
    Q = \Delta U + A, \label{eq:first_law_45}
\end{equation}
где:
\begin{itemize}
    \item $Q$ --- \textbf{количество теплоты}, сообщаемое системе (или отводимое от неё);
    \item $\Delta U = U_2 - U_1$ --- изменение \textbf{внутренней энергии} системы;
    \item $A$ --- \textbf{работа}, совершаемая системой над внешними телами.
\end{itemize}
\end{definition}

\begin{definition}
\textbf{Знаки величин в первом начале термодинамики:}
\begin{itemize}
    \item $Q > 0$ --- система \textbf{получает} теплоту (нагревается);
    \item $Q < 0$ --- система \textbf{отдаёт} теплоту (охлаждается);
    \item $A > 0$ --- система \textbf{совершает работу} над внешними телами (расширяется, $\Delta V > 0$);
    \item $A < 0$ --- внешние силы \textbf{совершают работу} над системой (сжатие, $\Delta V < 0$);
    \item $\Delta U > 0$ --- внутренняя энергия \textbf{увеличивается} (температура растёт);
    \item $\Delta U < 0$ --- внутренняя энергия \textbf{уменьшается} (температура падает).
\end{itemize}
\end{definition}

\begin{definition}
\textbf{Количеством теплоты} $Q$ называется энергия, передаваемая от одного тела к другому в процессе теплообмена (теплопроводность, конвекция, излучение) без совершения макроскопической работы. В отличие от внутренней энергии, теплота --- это \textbf{процессная величина}: она характеризует способ передачи энергии, а не состояние системы. Единица измерения --- джоуль (Дж).
\end{definition}

\begin{theorem}[Первое начало для изопроцессов]
\textbf{1. Изохорный процесс ($V = \text{const}$):}
\begin{equation}
    A = \int p \, dV = 0 \quad \Rightarrow \quad Q_V = \Delta U. \label{eq:isochoric_45}
\end{equation}
Вся подводимая теплота идёт на изменение внутренней энергии.

\textbf{2. Изобарный процесс ($p = \text{const}$):}
\begin{equation}
    A = p \Delta V = \nu R \Delta T \quad \Rightarrow \quad Q_p = \Delta U + p \Delta V. \label{eq:isobaric_45}
\end{equation}
Теплота расходуется на изменение внутренней энергии и совершение работы расширения.

\textbf{3. Изотермический процесс ($T = \text{const}$):}
\begin{equation}
    \Delta U = 0 \quad \Rightarrow \quad Q_T = A = \nu R T \ln \frac{V_2}{V_1}. \label{eq:isothermal_45}
\end{equation}
Вся подводимая теплота полностью превращается в механическую работу.

\textbf{4. Адиабатный процесс ($Q = 0$):}
\begin{equation}
    0 = \Delta U + A \quad \Rightarrow \quad A = -\Delta U. \label{eq:adiabatic_45}
\end{equation}
Работа совершается за счёт убыли внутренней энергии (температура падает при расширении).
\end{theorem}

\begin{remark}
Первое начало термодинамики запрещает создание \textbf{вечного двигателя первого рода} --- машины, которая совершала бы работу без затрат энергии извне или производила бы больше работы, чем полученная энергия. Математически это следует из \eqref{eq:first_law_45}: если $Q = 0$ и $\Delta U = 0$, то $A = 0$.
\end{remark}
